\section{A common potential argument}

We use the assumptions and notation of Section~1. Put
$c=\min_{|u|=1}\rho(u)>0$, $\Lambda=\max_{|u|=1}\rho(u)$,
and $C_x=x+[-1/2,1/2)^d$. Write $q_x$ for the conditional law
used at the first departure from $x$; this law can be random,
but its prescribed mean is fixed. Every estimate below is pathwise
uniform over its possible values. Set $\Delta=P-I$ for the
simple-random-walk generator and
$\bar\nabla_i f(x)=(f(x+e_i)-f(x-e_i))/2$.
For a bounded measurable $D$, define
\begin{equation}\label{eq:core-potential}
 U_D(y)=\frac{2\kappa}{\omega_d}
 \int_D\nabla\rho(v)\cdot\frac{v-y}{|v-y|^d}\,dv.
\end{equation}

\begin{lemma}[Cone, layers, and contact]\label{lem:core-geometry}
For $b\geq0$,
$U_{bK}(y)=2d\kappa(b-\rho(y))_+$.
If $E\subset\{b\leq\rho\leq b+a\}$ then
$\|U_E\|_\infty\leq2d\kappa a$; for arbitrary $E$,
$\|U_E\|_\infty\leq C|E|^{1/d}$.
Suppose $bK\subset D\subset B(0,R)$ up to null sets, $y_0\in\partial(bK)$ lies in
$\overline{D^c}$, and $E=D\setminus bK$. Put
$H=U_D(y_0)$ and $I=\int_E(\rho-b)$.
Then
\begin{equation}\label{eq:core-contact}
 H\geq0,\qquad I\leq CR^d H,\qquad
 \|U_E\|_\infty\leq C I^{1/(d+1)}.
\end{equation}
\end{lemma}
\begin{proof}
For a compactly supported Lipschitz $f$, polar coordinates about $y$ give
$\int\nabla f(v)\cdot(v-y)|v-y|^{-d}\,dv=-d\omega_df(y)$.
Apply this to $f=(b-\rho)_+$. On each ray $y+tu$, convexity makes
$t\mapsto\rho(y+tu)$ decrease and then increase. On either part, its total
variation while $b\leq\rho\leq b+a$ is at most $a$; the two signed parts
therefore give an integral between $-a$ and $a$. Integrating over directions
proves the layer bound. Rearrangement of $|v-y|^{1-d}$ proves the volume
bound. Finally
$\nabla\rho(v)\cdot(v-y_0)\geq\rho(v)-\rho(y_0)=\rho(v)-b$
on $E$, proving the first two assertions in \eqref{eq:core-contact}.
Split $E$ at $\rho=b+a$: the layer contributes at most $Ca$ and the rest,
of volume at most $I/a$, contributes at most $C(I/a)^{1/d}$.
Take $a=I^{1/(d+1)}$; the case $I=0$ follows directly.
\end{proof}

\begin{lemma}[Discrete potential representation]\label{lem:core-green}
Let $g$ be the potential kernel in dimension two and minus the Green
function in dimensions at least three, normalized by
$\Delta g=\mathbf1_{\{0\}}$, where $\Delta=P-I$ for simple random walk.
Write $g_y(x)=g(x-y)$ and $M^y=M^{g_y}$ for its Dynkin martingale;
for finite combinations of these kernels, extend $M^f$ linearly.
These martingales have bounded increments. Whenever
$\max_{j\leq t}|X_j|\leq R$, $R\geq2$, and $y\in\Z^d$ satisfies $|y|\leq R^{d+2}$,
\begin{equation}\label{eq:core-green}
 \ell_t(y)=U_{D_t}(y)-M_t^y+O(\log R),\qquad
 \langle M^y\rangle_t\leq C W_t(y),\quad
 W_t(y)=\sum_x\ell_t(x)(1+|x-y|)^{2-2d}.
\end{equation}
Moreover, for $|z-y|\leq\sqrt d$,
$|U_{D_t}(z)-U_{D_t}(y)|\leq C\log R$.
\end{lemma}
\begin{proof}
Standard Green estimates \citep{LawlerLimic2010} give
\begin{align*}
 \bar\nabla g(z)&=\frac2{\omega_d}\frac z{|z|^d}+O((1+|z|)^{-d})
 \quad(z\ne0),\\
 |g(z+e)-g(z)|&\leq C(1+|z|)^{1-d},\\
 |g(z+e_i)+g(z-e_i)-2g(z)|&\leq C(1+|z|)^{-d}.
\end{align*}
Write $s_i(x)=q_x(e_i)+q_x(-e_i)$. Exactly,
\[
 (Q_x-P)g_y(x)=-\kappa\lambda(x)\xi(x)\cdot\bar\nabla g(x-y)
 +\frac12\sum_i(s_i(x)-1/d)
 [g(x+e_i-y)+g(x-e_i-y)-2g(x-y)].
\]
Dynkin's formula, grouping first departures by site, gives the first
identity in \eqref{eq:core-green} with the discrete drift potential.
The displayed second differences sum to $O(\log R)$. The endpoint
$g(X_t-y)-g(-y)$ has this bound too. By dyadic annuli,
\begin{equation}\label{eq:core-defect-potential}
 \sup_y\sum_{x\in B(0,CR)}
 |\lambda(x)-1|(1+|x-y|)^{1-d}\leq C\log R.
\end{equation}
Thus replacing $\lambda$ by one costs $O(\log R)$, even for an arbitrary
subset of visited defect sites. Replacing the Green gradient by its
leading term, and then by cell averages, has the same cost.

For completeness the subgradient replacement also costs only $O(\log R)$,
including at corners of $K$. The positive measure $\mu=\Delta_{\rm cont}\rho$
satisfies $\mu(B(y,s))\leq Cs^{d-1}$, by testing against a cutoff and using
$|\nabla\rho|\leq\Lambda$. The convex-cell estimate is
\begin{equation}\label{eq:core-cell}
 \int_{C_x}|\nabla\rho(v)-\xi(x)|\,dv
 \leq C\mu(B(x,6\sqrt d)).
\end{equation}
Indeed, $h(v)=\rho(v)-\rho(x)-\xi(x)\cdot(v-x)$ is nonnegative, convex,
and zero at $x$. Convexity along coordinate lines gives
$|\partial_i h(v)|\leq h(v+e_i)+h(v-e_i)$, so the left side is at most
$C\int_{B(x,2\sqrt d)}h$. Mollify $h$ and put $A=2\sqrt d$. Convexity gives
\[
 h_\varepsilon(x+tu)\le h_\varepsilon(x)
       +t\,\partial_r h_\varepsilon(x+2Au)
       \qquad(0\le t\le A,\ |u|=1).
\]
Integration in $t,u$ and the divergence theorem bound the integral by
$C h_\varepsilon(x)+C\mu(B(x,6\sqrt d))$.
Letting $\varepsilon\downarrow0$ proves \eqref{eq:core-cell}.
Bound the cells close to $y$ directly; bounded overlap and dyadic annuli
bound the remaining error by
$C\int_{1\leq|z-y|\leq CR^{d+2}}(1+|z-y|)^{1-d}\,d\mu(z)=O(\log R)$.
The increment estimate gives the bracket bound. Splitting the integral
at distance $2\sqrt d$ and differentiating the far kernel gives the
cell modulus.
\end{proof}

\begin{lemma}[Exit bootstrap]\label{lem:core-bootstrap}
Put $\tau_R=\inf\{t:\rho(X_t)\geq R\}$. For every $p>0$, outside an event
of probability $CR^{-p}$, simultaneously for $t\leq\tau_R$,
\begin{equation}\label{eq:core-bootstrap}
 \max_x\ell_t(x)\leq CR,\qquad
 \tau_R\leq CR^{d+1},\qquad
 \sup_{|y|\leq R^{d+2}}|M_t^y|\leq C\sqrt R\log R.
\end{equation}
Consequently
$\|\widetilde\ell_t-U_{D_t}\|_\infty\leq C\sqrt R\log R$ on
$B(0,CR)$.
\end{lemma}
\begin{proof}
To see that the exit time is finite, stop the compensated coordinate
martingale $Z_t^i=X_t^i+\kappa\sum_{x\in A_t}\lambda(x)\xi_i(x)$
at $\tau_R$. It is bounded by a deterministic constant $C_R$.
There are at most $m_R=O(R^d)$ first departures before exit, so its
bracket is at least $((t\wedge\tau_R)-m_R)/d$. Consequently
$\E(t\wedge\tau_R)\le m_R+dC_R^2$, and $\tau_R<\infty$ almost surely. Before exit, all departed sites lie in $B(0,CR)$.
The exact discrete Dynkin identity and the first Green increment bound
give $\ell_t(y)\leq |M_t^y|+CR$ without any preliminary clock estimate.
Stop also when one of the displayed martingales first reaches $R$ in
absolute value. Up to this stopping time, $\max_x\ell(x)\leq C'R$,
and hence the clock is at most $C'R^{d+1}$ and every bracket at most
$C'R\log R$. Freedman's maximal inequality \citep[Theorem~1.6]{Freedman1975} and a union bound over the
polynomially many poles give
$\sup|M^y|\leq C_p\sqrt R\log R$ with probability $1-CR^{-p}$.
This is smaller than $R$ for large $R$, so the additional stop is never
reached before $\tau_R$. The representation and cell modulus finish.
\end{proof}

\begin{lemma}[Exponential crossing and projection]\label{lem:core-crossing}
For every $p>0$ and integer $N\geq2$, outside an event of probability
$CN^{-p}$, simultaneously for $t\leq N$ and integers $0\leq k\leq C_0N$,
\begin{equation}\label{eq:core-crossing}
 \max_{j\leq t}\operatorname{dist}(X_j,kK)
 \leq C\left(1+\max_{\rho(x)\geq k}\ell_t(x)\right)\log N.
\end{equation}
Here $C_0$ is any fixed constant.
\end{lemma}
\begin{proof}
We first record the short crossing estimate. Fix a unit vector $u$, $a\geq0$,
and suppose all departed sites in the cap $u\cdot x>a$ have
$u\cdot\lambda(x)\xi(x)\geq\alpha>0$. With
$f(x)=e^{\theta u\cdot x}$, conditional increments satisfy
\[
 \mathbb E[f(X_{j+1})-f(X_j)\mid\mathcal F_j]
 \leq f(X_j)(C\theta^2-\kappa\alpha\theta I_j)
 \quad (u\cdot X_j>a),
\]
where $I_j$ indicates a first departure. Stop at time $N$, on entering
a cap site violating the drift condition, when a cap local time exceeds
$L$, or when $u\cdot X$ reaches $a+h$. Choose $\theta=c/(L+1)$.
Summing by site makes the total contribution from the cap nonpositive:
each departed site supplies its cookie once and at most $L+1$ visits.
The other at most $N$ increments contribute at most
$CN e^{\theta(a+1)}$. Optional stopping of the Dynkin martingale gives
\begin{equation}\label{eq:core-exp-crossing}
 \mathbb P\left(\begin{array}{c}
  \text{the level }a+h\text{ is reached by }N,\text{ all previously visited}\\
  \text{cap sites satisfy the drift condition, and cap local times}\leq L
 \end{array}\right)
 \leq CN\exp\{-ch/(L+1)\}.
\end{equation}
The stopping argument does not condition on the event in this display.

Let $w=\max_{j\leq t}\operatorname{dist}(X_j,kK)$ and choose $x_*=X_j$
attaining it. When $w>0$, let $\pi$ denote Euclidean projection onto $kK$,
$u=(x_*-\pi(x_*))/w$, and $\beta=u\cdot\pi(x_*)$.
Then $u\cdot X_j\leq T:=\beta+w$ for all $j\leq t$ and $\beta\geq0$.
Put $\eta=c^2/(8\Lambda^2)$. If $x$ is a site visited by time $t$
and $u\cdot x>T-\eta w$, write $v=(x-\pi(x))/|x-\pi(x)|$.
The projection inequalities give $u\cdot v>1-\eta$; moreover
\[
 \xi(x)\cdot v\geq
 \frac{\rho(x)-k}{\operatorname{dist}(x,kK)}\geq c,
\]
where the last inequality follows by comparing $\pi(x)$ with
$kx/\rho(x)$. Thus $u\cdot\lambda(x)\xi(x)\geq\lambda_-c/2$.
All these cap sites satisfy $\rho(x)>k$.
Use \eqref{eq:core-exp-crossing} with $a$ the least integer at least
$T-\eta w$, and with $L$ the integer above the largest relevant local time.
The directions $u$ belong to the deterministic list formed by projecting
the lattice sites $|x|\leq N$ onto $kK$, for integer $k\leq C_0N$.
This list, the integer levels $a\leq CN$, and $1\leq L\leq N+1$ have
polynomial total size. For each triple choose the deterministic length
$C_p(L+1)\log N$ in \eqref{eq:core-exp-crossing}. A union bound excludes
all these crossings. The drift condition is needed only at visited
sites, since the stopped estimate ends on entering any violating site. Since $T-a\geq\eta w-1$, this proves
\eqref{eq:core-crossing}; bounded $w$ is absorbed into the constant.
\end{proof}

\begin{proposition}[Coarse scale, with polynomial tails]\label{prop:core-coarse}
For every $p>0$ there are $c_p,C_p>0$ such that, with probability at least
$1-C_pR^{-p}$,
\begin{equation}\label{eq:core-clock}
 c_pR^{d+1}\leq\tau_R\leq C_pR^{d+1}.
\end{equation}
Consequently, with $r_n=((d+1)n/(2d\kappa|K|))^{1/(d+1)}$, with probability
at least $1-C_p n^{-p}$,
\begin{equation}\label{eq:core-coarse}
 \max_{j\leq n}|X_j|\leq C_pr_n,\quad
 c_pr_n\leq\max_x\ell_n(x)\leq C_pr_n,\quad
 c_pr_n^d\leq|A_n|\leq C_pr_n^d.
\end{equation}
\end{proposition}
\begin{proof}
Work on the events of the preceding two lemmas with a deterministic
horizon $N=\lceil CR^{d+1}\rceil$. At time $\tau_R$, let
$b=\inf_{y\notin D_{\tau_R}}\rho(y)$, and choose a contact point
$y_0\in\partial(bK)\cap\overline{D_{\tau_R}^c}$ in the closure of an
unoccupied cell. Continuity and the cell estimate imply
$0\leq H=U_{D_{\tau_R}}(y_0)\leq\delta_R=C\sqrt R\log R$.
Lemma~\ref{lem:core-geometry} therefore gives
\[
 \|\widetilde\ell_{\tau_R}-2d\kappa(b-\rho)_+\|_\infty
 \leq\delta_R+C(R^d\delta_R)^{1/(d+1)}
 =:\varepsilon_R=o(R/\log R).
\]
Use \eqref{eq:core-crossing} with $k=\lfloor b\rfloor$; outside $kK$ the
local times are at most $C+\varepsilon_R$. Since
$\rho(X_{\tau_R})\geq R$ and $\rho$ is $\Lambda$-Lipschitz,
\[
 R-b\leq C+C(\varepsilon_R+1)\log R=o(R).
\]
The cone is supported on $\{\rho<b\}\subset D_{\tau_R}$.
Integrating its approximation over $D_{\tau_R}\subset B(0,CR)$ gives
\[
 \tau_R=\frac{2d\kappa|K|}{d+1}b^{d+1}
 +O(R^d\varepsilon_R)\asymp R^{d+1},
\]
proving \eqref{eq:core-clock}. To obtain \eqref{eq:core-coarse}, choose
$R=\lceil C n^{1/(d+1)}\rceil$ so that the lower clock bound exceeds $n$,
and apply the exit event with exponent $(d+1)(p+1)$. Then $n<\tau_R$,
so the spatial and local-time upper bounds follow. The remaining bounds
follow from $n=\sum_x\ell_n(x)\leq|A_n|\max_x\ell_n(x)$ and
$|A_n|\leq C(1+\max_{j\leq n}|X_j|)^d$.
\end{proof}
